\chapter{Extensions: urban commuting, terminals, and Seattle}
\label{ch:extensions}

\section{One transport block, different spatial closures}

The package shares route, modal, congestion, and primitive pass-through
operators across economic geography and urban commuting. The models differ in
their equilibrium residuals, state variables, and welfare projection.

The urban state is
\[
 z_U=(d\log l^R,d\log l^F,d\log\chi),
 \qquad
 d\log W=-\frac{d\log\chi}{\theta}.
\]
Residence mass \(l^R\) and workplace mass \(l^F\) can differ by location.
The route kernel therefore reconstructs origin-destination commuting with
distinct source and destination margins.

The urban transport equations remain:
\[
 \kappa_e=\left(\sum_m\kappa_{e,m}^{-\eta}\right)^{-1/\eta},
\qquad
 \tau_{ij}^{-\theta}
 =\mathbf 1\{i=j\}+\sum_k\kappa_{ik}^{-\theta}\tau_{kj}^{-\theta}.
\]
The shared Schur complement eliminates route, mode, and congestion quantities,
while the urban Jacobian \(J_U^0\) and welfare row \(q_U\) replace the trade
model's spatial block.

\section{Urban verification}

The one-mode regression oracle is a separate implementation of the
Allen--Arkolakis urban derivative \citep{AA_2022restud}. The shared engine must
reproduce its state response and welfare derivative below \(10^{-10}\),
including the familiar \(1+\theta\lambda\) congestion denominator.

The multimodal urban example then compares the analytical IFT with an
independently solved nonlinear system under \code{NC}, \code{NT}, \code{F},
\code{FM}, and \code{FR}. Tests cover:
\begin{itemize}
  \item distinct residence and workplace margins;
  \item road and transit modes;
  \item edge and endpoint-terminal congestion;
  \item one-mode, unique-route, no-congestion, and zero-terminal limits;
  \item node and mode permutations.
\end{itemize}

Run the self-contained example with:
\begin{lstlisting}[language=bash]
julia --project=. bin/tnw.jl validate \
  examples/urban_multimodal/config.toml
julia --project=. bin/tnw.jl decompose \
  examples/urban_multimodal/config.toml
\end{lstlisting}

\section{Endpoint-terminal congestion}

Edge congestion assigns traffic feedback to a particular edge-mode. Terminal
congestion aggregates affected flows using a terminal and applies the cost
response to incident edge-mode pairs. The terminal identifiers are explicit
input fields; geographic co-location is not enough.

In the RSUE extension, rail traffic loads the terminals at both endpoints with
elasticity \(0.096\). The main paper result keeps this channel off. Sensitivity
to terminal congestion is therefore an extension diagnostic, not a headline
parameter sensitivity.

\section{Seattle multimodal candidate}

The Seattle workflow combines:
\begin{itemize}
  \item the published 217-location Allen--Arkolakis urban geography;
  \item 2017 LODES residence-workplace commuting;
  \item 2017 ACS residential commute-mode shares;
  \item a hash-pinned June 2017 King County GTFS feed;
  \item a Metro route report used as external comparison.
\end{itemize}

The baseline uses \(\eta=1.099\), transferred from the freight application.
It is not estimated from Seattle data. Baseline modal shifters reproduce the
constructed shares; \(\eta\) controls the local response after an improvement.
The workflow therefore reports sensitivity at \(0.75\), \(0.90\), \(1.099\),
\(1.25\), and \(1.40\).

GTFS describes stops, trips, shapes, and schedules. It does not provide
passenger counts. ACS supplies residential mode shares, not route-level
ridership. The builder routes a common commuter population to obtain a
model-consistent road and transit baseline, records fallback and coverage, and
uses the Metro table only as an external comparison.

\section{Transit policies}

The workflow evaluates bus, Link light rail and streetcar, and ferry policies
separately and together. A link policy lowers primitive generalized cost on
one directed transit edge-mode. Reciprocal pairs are aggregated only when both
directions exist.

A named route corridor is a policy bundle:
\[
 E_b=\sum_{(e,m)\in b}w_{e,m}E_{e,m}.
\]
Shared segments change the aggregate mode cost on that segment. The result is
not interpreted as a shock exclusive to one operator's vehicles.

Maps distinguish road, bus, Link light rail, Seattle Streetcar, and ferry by
color. Line width represents the extended welfare gain, not ridership.
Subway or light-rail vision maps that include proposed future lines are not a
baseline description of the 2017 network.

\section{What remains empirical work}

The package supplies a verified multimodal urban derivative. A publication
quality Seattle transit application would still benefit from:
\begin{itemize}
  \item observed passenger counts by route or segment;
  \item mode-specific generalized costs and service quality;
  \item a Seattle-specific estimate or external calibration of \(\eta\);
  \item crowding or terminal-congestion elasticities;
  \item validation against untargeted commuting and ridership outcomes.
\end{itemize}
These needs concern calibration and measurement, not a relabeling of the IFT.
